\chapter{Arquitectura Global del Sistema}
Este capítulo actúa como el pegamento de la tesis, demostrando la visión de ingeniería de sistemas y redes. Justificaremos el abandono del paradigma "Cloud-Centric" a favor de una topología híbrida orientada al \textit{Edge}.

\section{Vista General}

Concebimos la arquitectura global del sistema propuesto como una solución híbrida y jerárquica que rompe deliberadamente con el paradigma tradicional estrictamente centralizado en la nube, desplazando el núcleo de la lógica de cómputo e inferencia hacia el \textit{Edge Computing}. Diseñamos esta aproximación específicamente para responder a las restricciones operativas del entorno agrícola español, el cual está caracterizado por una infraestructura de conectividad intermitente en redes \gls{wan}, recursos energéticos limitados en campo y la necesidad imperante de mantener el control absoluto sobre la información generada, erradicando así el pago de tarifas recurrentes a proveedores externos.

Para articular este flujo y garantizar una separación clara de responsabilidades computacionales, hemos estructurado la topología de red y procesamiento rigurosamente en tres capas físicas y lógicas interconectadas:

\begin{enumerate}
    \item \textbf{Capa de Instrumentación y Adquisición Física} | Nodos IoT: Esta capa constituye el nivel más periférico, denominado \textit{Deep Edge}, y está materializada por estaciones agro-climáticas portadoras del \textit{firmware} \textit{Savia}. Operamos estas estaciones mediante microcontroladores de arquitectura avanzada, empleando en este caso la plataforma \textit{Raspberry Pi Pico 2 W} con el sistema en chip RP2350. Buscamos con el diseño de estos nodos la máxima eficiencia energética y resiliencia, gestionando la escasez de recursos mediante estados de suspensión profunda intermitentes y ciclos de trabajo optimizados.
    
    A nivel funcional, el \textit{firmware} embebido \textit{Savia}, desarrollado en C sobre el SDK nativo por nuestro compañero de proyecto César Almeida, opera de manera determinista asumiendo cuatro responsabilidades concurrentes:
    \begin{itemize}
        \item \textbf{Adquisición y Sensorización:} Establecemos interfaz directa con instrumentación agrícola, como sondas de humedad y temperatura del suelo, utilizando estándares de instrumentación multipunto como el protocolo SDI-12.
        \item \textbf{Agregación y Persistencia Local:} Preprocesamos las series temporales brutas para la consolidación de medias horarias. Custodiamos estos datos temporalmente en estructuras de búfer circular sobre memoria no volátil, garantizando la preservación de la telemetría durante un mínimo de 48 horas ante caídas prolongadas de los enlaces de radio.
        \item \textbf{Inferencia sobre el dispositivo:} Aprovechamos las capacidades de cómputo del microcontrolador, incluyendo su unidad de punto flotante y la disponibilidad de 520 kilobytes de memoria RAM, para ejecutar localmente modelos de aprendizaje profundo. Específicamente, utilizamos redes neuronales recurrentes \gls{lstm} cuantizadas \footnote{\gls{lstm} Cuantizado: El mismo modelo cuantizado, sin pérdidas que afecten a la confiabilidad final, lo utilizamos más adelante para la inferencia en la aplicación de usuario.} mediante \textit{TensorFlow Lite for Microcontrollers}. Esto dota a la estación de autonomía analítica para generar pronósticos a 24 horas sobre el estado hídrico sin depender del acceso a la nube.
        \item \textbf{Interfaces de Comunicación:} Exponemos la telemetría histórica y los vectores de predicción hacia la capa intermedia a través de perfiles \gls{ble} GATT. Adicionalmente, contamos con un canal de red de área amplia y baja potencia, basado en tecnología LoRa, para la emisión periódica de tramas de estado crítico siempre y cuando exista una capacidad de señal viable para el transceptor.
    \end{itemize}
    
    \item \textbf{Capa Logística y de Inferencia} | \textit{Edge}: Materializamos esta capa a través del teléfono inteligente del agricultor, dispositivo que ejecuta una aplicación nativa de alto rendimiento desarrollada sobre el entorno \textit{Flutter/Dart}. En nuestra topología, concebimos esta capa como el cuartel general logístico y motor principal de ejecución del sistema. A diferencia del diseño clásico donde el terminal móvil asume únicamente el papel de interfaz gráfica pasiva orientada al usuario final, en nuestra arquitectura le otorgamos un rol activo y crítico como nodo \textit{Edge}. De este modo, concentramos la lógica operativa y distribuimos la carga a través de las siguientes funciones concurrentes:
    \begin{itemize}
        \item \textbf{Ingestión de Datos en Proximidad:} Gestionamos de forma proactiva la comunicación de corto alcance con los nodos de campo o capa inferior. La aplicación asume la responsabilidad de ingestar la telemetría binaria y el registro histórico acumulado a través de enlaces asíncronos \gls{ble}. Así, garantizamos la recolección de los datos directamente en la parcela, sin requerir cobertura celular en las zonas de cultivo.
        \item \textbf{Fusión de Datos y Contexto Climático:} Enriquecemos la información física de las sondas con el contexto macroclimático. Habilitamos la conexión síncrona de la aplicación con la plataforma Open-Meteo vía red celular \gls{wan} con el objetivo de extraer predicciones meteorológicas. De esta manera, fusionamos localmente las variables microclimáticas en tiempo real con las previsiones a corto plazo.
        \item \textbf{Inferencia Predictiva en el \textit{Edge}:} Ejecutamos el flujo completo de aprendizaje automático en la propia memoria del dispositivo mediante un ecosistema de inferencia dual. Por un lado, integramos un motor embebido basado en \textit{TensorFlow Lite}, capacitado para interpretar modelos \gls{lstm} cuantizados. Por otro lado, incorporamos la capacidad de evaluar modelos \textit{Random Forest}. Logramos esto gracias a que exportamos previamente estos modelos a descripciones estructuradas en formato JSON, lo cual nos permite transformarlos y reconstruirlos como algoritmos dinámicos en pleno tiempo de ejecución. Esta combinación tecnológica dota a la aplicación de la habilidad para evaluar el conjunto de datos fusionado y generar un veredicto predictivo de riego en cuestión de milisegundos. Aunque ejecutamos la computación íntegramente en el dispositivo local, este proceso de inferencia constituye el único eslabón operativo que requiere conectividad celular en tiempo real para obtener las predicciones climáticas, mientras que delegamos tareas como la descarga \gls{ble} o la sincronización general con la base de datos a rutinas asíncronas totalmente independientes de la cobertura inmediata.
        \item \textbf{Pasarela de Sincronización:} Utilizamos la aplicación como puente lógico bidireccional hacia la capa superior de \textit{Fog Computing} conformada por el servidor auto-alojable. Al detectar una conexión estable a Internet, orquestamos la transmisión asíncrona de las series temporales recolectadas hacia la base de datos relacional centralizada, habilitando de esta forma la red de cooperación entre agricultores. Simultáneamente, consultamos y descargamos de manera transparente posibles actualizaciones en las lógicas de nuestros modelos predictivos.
    \end{itemize}
    
    \item \textbf{Capa de Consolidación y Sincronización Global | Servidor Central Auto-Alojable:} Constituye el eslabón de agregación a nivel de finca o cooperativa, actuando como \textit{Fog Computing}. Materializamos esta capa bajo una filosofía de soberanía de datos que exime al agricultor de la dependencia técnica y los costes recurrentes asociados a plataformas comerciales en la nube. Este componente consiste en un servidor telemático ligero y auto-alojable, el cual hemos desarrollado íntegramente sobre el entorno de ejecución Node.js. 

    Para mantener una huella computacional mínima y atenernos estrictamente a principios de ingeniería ágil como el paradigma \textit{YAGNI}, es decir, \textit{You Aren't Gonna Need It}, hemos prescindido de \textit{frameworks} monolíticos de terceros en la arquitectura. En su lugar, resolvemos el enrutamiento y procesamiento concurrente \gls{http} directamente sobre las bibliotecas estándar de red y sistema de archivos. Dotamos a este servidor de tres funciones fundamentales:
    \begin{itemize}
        \item \textbf{Persistencia y Cooperación Distribuida:} Alojamos un motor de base de datos relacional transaccional sin configuración adicional basado en SQLite, el cual actúa como repositorio maestro unificado. Esta base de datos facilita la sincronización asíncrona de las series temporales ingeridas por los nodos móviles, habilitando la consolidación de datos y posibilitando redes de inteligencia compartida entre agricultores sin ceder la propiedad de la información.
        \item \textbf{Distribución de Modelos Predictivos:} Utilizamos el servidor como repositorio central para la gestión y distribución de modelos de \textit{Machine Learning}, albergando estructuras como árboles de decisión \textit{Random Forest} convertidos a formato JSON. Gracias a esta arquitectura, permitimos a los dispositivos de la capa \textit{Edge} consultar, descargar y actualizar dinámicamente sus lógicas de predicción estacionales de forma transparente.
        \item \textbf{Seguridad Perimetral y Despliegue Agnóstico:} Canalizamos todo el intercambio de información a través de una \gls{api} RESTful segura. Hemos implementado en ella una autenticación estricta mediante claves precompartidas o \textit{API Keys}, mitigación de ataques de denegación de servicio mediante limitación de tasa de peticiones en memoria, y mecanismos rigurosos contra la escalada de privilegios y saltos de directorio. A nivel logístico, encapsulamos el servicio empleando tecnología de contenerización mediante Docker, con lo que garantizamos despliegues resilientes, portables y aislados bajo perfiles de usuario sin privilegios.
    \end{itemize}
\end{enumerate}

Esta distribución estratégica de la carga de trabajo no solo optimiza drásticamente el uso del ancho de banda y la energía de los dispositivos empotrados, sino que dota al sistema de una alta tolerancia a fallos. La pérdida de conectividad con la tercera capa o servidor central no impide a la segunda capa móvil descargar datos de campo vía \gls{ble} ni inferir localmente las necesidades hídricas, siempre y cuando disponga de acceso celular básico para consultar los servicios meteorológicos. De este modo, la arquitectura física mapea perfectamente los requerimientos de resiliencia y alto rendimiento exigidos en la Agricultura 4.0 contemporánea.

\section{Logística}
En el contexto de nuestra arquitectura híbrida, definimos el rol del teléfono inteligente no como un simple visor pasivo de métricas, sino que lo concebimos genuinamente como el Cuartel General Logístico de aprendizaje automático en el \textit{Edge}. Atribuimos a este dispositivo móvil la responsabilidad de orquestar el descubrimiento de la red de sensores dispersos en el campo. Mediante rutinas de escaneo en segundo plano, la aplicación identifica y negocia la conexión con cada estación agro-climática disponible, gestionando de manera eficiente la ingesta de telemetría intermitente. Una vez recolectados los datos, el terminal móvil asume la carga computacional para la ejecución de la inteligencia artificial, evaluando las variables y emitiendo directrices operativas sin depender de una infraestructura externa. Esta centralización lógica en el borde de la red nos permite minimizar los cuellos de botella y asegurar la viabilidad del sistema en zonas rurales aisladas.

\section{Interacción entre Componentes}
Mapeamos en detalle las interfaces de comunicación para comprender el flujo de información a través del sistema. Detallamos rigurosamente qué protocolos unen cada nodo operativo para garantizar una interoperabilidad robusta. En primer lugar, empleamos el estándar \gls{ble} para vincular la capa física de instrumentación con el nodo \textit{Edge}. Este protocolo de muy bajo consumo energético resulta idóneo para la descarga asíncrona de datos en proximidad. En segundo lugar, integramos peticiones basadas en \gls{http} bajo un diseño RESTful para conectar el dispositivo \textit{Edge} con las interfaces de programación de aplicaciones meteorológicas de terceros. Por último, utilizamos la suite de protocolos TCP/IP sobre redes celulares estandarizadas para enlazar el nodo móvil con el servidor central auto-alojable, facilitando así la sincronización global y la distribución de los modelos predictivos.

\section{Flujo de Datos}
Trazamos el ciclo de vida completo de un paquete de datos para ilustrar el cauce o \textit{pipeline} de procesamiento de nuestra plataforma. El recorrido comienza en el momento en el que la estación IoT adquiere y filtra las lecturas de los sensores físicos. A continuación, el microcontrolador comprime esta información en un formato de empaquetado binario altamente optimizado para minimizar la sobrecarga de red. Posteriormente, ejecutamos la transmisión asíncrona de estas tramas hacia el dispositivo móvil mediante radiofrecuencia de corto alcance. Una vez que el teléfono inteligente recepciona los paquetes, procedemos a su almacenamiento temporal en bases de datos locales, asegurando así la persistencia de los registros. A partir de este repositorio intermedio, el motor de aprendizaje automático ingiere las series de datos para emitir sus inferencias predictivas. Finalmente, al establecer un enlace a Internet estable, el sistema coordina la subida asíncrona de todo el histórico hacia el servidor central, cerrando con éxito el ciclo de consolidación de la información.

\section{Seguridad y Privacidad de Datos}
Analizamos exhaustivamente la superficie de ataque de nuestra infraestructura para garantizar la integridad de las operaciones agrícolas. Abordamos la seguridad desde el primer eslabón de la cadena mediante la configuración estricta de la transmisión \gls{ble}. En esta etapa, implementamos mecanismos criptográficos de emparejamiento validado que previenen la intercepción maliciosa o la suplantación de identidad de los sensores de campo. A nivel de infraestructura superior, protegemos los accesos al servidor central exigiendo una autenticación robusta en la interfaz de programación REST mediante el uso de \textit{JSON Web Tokens} o claves precompartidas de alta entropía. Más allá de la protección perimetral, destacamos la ventaja inherente de privacidad que proporciona el paradigma \textit{Edge}. Al procesar los datos crudos directamente en el dispositivo local, evitamos exponer la lógica de negocio y los hábitos de cultivo en servidores comerciales en la nube, garantizando que los agricultores conserven en todo momento la soberanía absoluta e inalienable sobre su propia información.
